% ================================================================================================ capítulo 12
\chapter{El satélite GOES y el seeing desde el espacio}\label{cap:satelite_teoria}

\begin{resumen}
\textbf{En este capítulo:} qué mide el instrumento ABI de los satélites GOES, qué es una temperatura de brillo, cómo se
pasa de una latitud y una longitud a un píxel de la imagen, cómo lee TurbuLENS solo un pedazo de un archivo remoto, y
el método de Cavazzani, Ortolani y Zitelli para estimar el seeing con tres canales infrarrojos.
\end{resumen}

\section{GOES-Este y su instrumento ABI}

Los satélites GOES (\emph{Geostationary Operational Environmental Satellites}) de la NOAA giran sobre el ecuador a
35\,786~km de altura, con la misma velocidad angular que la Tierra: desde el suelo parecen quietos. El que mira a América
desde unos 75° de longitud oeste se llama \textbf{GOES-Este}: fue GOES-16 desde diciembre de 2017 hasta el 4 de abril
de 2025, y desde entonces es GOES-19.

Su cámara principal, el ABI (\emph{Advanced Baseline Imager}), toma una imagen del \textbf{disco completo} de la Tierra cada
10 minutos en \textbf{16 canales}: dos visibles, cuatro del infrarrojo cercano y diez del infrarrojo térmico. En el
infrarrojo térmico cada píxel mide unos 2~km bajo el satélite, y un poco más hacia el borde del disco; Bogotá está muy
cerca del punto bajo el satélite, así que allí los píxeles son de casi 2~km. TurbuLENS
usa el producto \codigo{ABI-L2-CMIPF}, las imágenes ya calibradas del disco completo:
\begin{itemize}
\item en los canales 1 a 6 el valor es la \textbf{reflectancia} (la fracción de la luz del Sol que se refleja; de noche no
sirve);
\item en los canales 7 a 16, la \textbf{temperatura de brillo} en kelvin, de día y de noche.
\end{itemize}

\subsection{Temperatura de brillo}

La temperatura de brillo de un píxel es la temperatura que tendría un cuerpo negro que emitiera, en la longitud de onda
del canal, la misma radiación que mide el satélite. Si lo que ve el canal es la superficie despejada, es parecida a la
temperatura del suelo; si es el tope de una nube alta, es mucho más fría. En el canal 13 (10,3~µm, el «infrarrojo
limpio»), el aire es casi transparente y se ve hasta la superficie o hasta el tope de las nubes:

\begin{table}[H]
\centering\small
\begin{tabular}{@{}l l@{}}
\toprule
Temperatura en el canal 13 & Lectura\\
\midrule
menos de $-40$~°C & Topes de nubes altas y frías\\
de $-40$ a $-10$~°C & Nubes medias o altas\\
de $-10$ a $5$~°C & Posibles nubes bajas o neblina\\
más de $5$~°C & Probablemente despejado (se ve la superficie)\\
\bottomrule
\end{tabular}
\caption{La lectura sencilla que muestra TurbuLENS para los canales infrarrojos limpios.}
\end{table}

En otros canales el aire absorbe: el canal 9 (6,9~µm) ve el vapor de agua de la troposfera media, así que su temperatura
es la de esa capa y no la del suelo; el canal 16 (13,3~µm) está en una banda del CO$_2$ y mezcla la emisión de capas bajas
y medias. Cada canal, en la práctica, «mira» a una altura distinta, y eso es lo que aprovecha el método del seeing.

\section{Del lugar al píxel}

Las imágenes del ABI no están en latitud y longitud sino en \textbf{ángulos de barrido} $(x, y)$: hacia dónde apuntaba el
espejo del instrumento. La guía de usuario de GOES-R da la conversión para un elipsoide terrestre (GRS80, con radio
ecuatorial $r_\text{eq}$ y polar $r_\text{pol}$) y un satélite a una distancia $H$ del centro de la Tierra, sobre la
longitud $\lambda_0$. Para un punto de latitud $\varphi$ y longitud $\lambda$:
\begin{gather}
\varphi_c = \arctan\!\left(\frac{r_\text{pol}^2}{r_\text{eq}^2}\tan\varphi\right),\qquad
r_c = \frac{r_\text{pol}}{\sqrt{1 - e^2\cos^2\varphi_c}},\qquad e^2 = \frac{r_\text{eq}^2 - r_\text{pol}^2}{r_\text{eq}^2},\\
s_x = H - r_c\cos\varphi_c\cos(\lambda-\lambda_0),\qquad s_y = -r_c\cos\varphi_c\sin(\lambda-\lambda_0),\qquad
s_z = r_c\sin\varphi_c,\\
y = \arctan\frac{s_z}{s_x},\qquad x = \arcsin\frac{-s_y}{\sqrt{s_x^2 + s_y^2 + s_z^2}} .
\end{gather}
El punto se ve desde el satélite si $H(H - s_x) \ge s_y^2 + (r_\text{eq}^2/r_\text{pol}^2)\,s_z^2$. La columna y la fila
salen de los ángulos con el primer valor y el paso de los ejes del archivo: $\text{columna} = (x - x_0)/\Delta x$. Todos
los parámetros ($H$, los radios, $\lambda_0$, los ejes) se leen de cada archivo. Con las mismas fórmulas, aplicadas a miles
de puntos a la vez, TurbuLENS dibuja sobre el recorte las costas, las fronteras y los departamentos de Natural Earth.

\section{Leer solo lo necesario}\label{sec:lectura_parcial}

Los archivos del ABI son NetCDF-4, que por dentro es HDF5: una cabecera y la imagen partida en bloques comprimidos. Para
leer un recorte de 600~km no hace falta el archivo entero, solo la cabecera y los bloques que lo cubren. El servidor de
la NOAA (en Amazon S3) acepta pedidos con la cabecera HTTP \codigo{Range: bytes=a-b}, que devuelve solo ese pedazo.

TurbuLENS presenta el archivo remoto como si fuera un archivo local: cuando el lector de HDF5 (pyfive, escrito en Python
puro) pide leer en una posición, se descarga el bloque de 128~KB que la contiene, y se guarda en memoria por si se vuelve a
necesitar (sección~\ref{sec:codigo_rango}). Así una imagen de unos 25~MB se lee bajando unos 2~MB. Si el servidor no
acepta pedidos parciales, TurbuLENS descarga el archivo completo.

\section{El seeing desde el satélite}\label{sec:metodo_cavazzani}

Cavazzani, Ortolani y Zitelli (2012) propusieron estimar el seeing con las imágenes infrarrojas nocturnas de un satélite
geoestacionario, y lo compararon con el seeing medido en Cerro Paranal (Chile) y en el Roque de los Muchachos (La Palma),
con una correlación cercana al 90\,\%. La idea es que tres canales que ven a alturas distintas permiten seguir la
estabilidad de la columna de aire.

\subsection{La función de estabilidad}

Con las temperaturas de brillo de los canales 9, 13 y 16, promediadas en un cuadro alrededor del sitio (1°~×~1° en el
artículo), se forma
\begin{equation}
F = T_9 - (T_{16} - T_{13}) .
\end{equation}
Durante la noche, $F$ cambia lentamente con el enfriamiento de la atmósfera. Lo que indica turbulencia son las
\textbf{desviaciones} de ese cambio lento: con todas las horas válidas de la noche se ajusta una recta $F_\text{tend}(t)$
por mínimos cuadrados, y en cada hora se toma $|F - F_\text{tend}|$. TurbuLENS necesita al menos tres horas para la recta;
con una o dos usa su promedio.

\subsection{De \texorpdfstring{$F$}{F} al seeing}

El método relaciona la desviación con la turbulencia integrada entre 4 y 8~km de altura:
\begin{equation}
\int_{4\,\text{km}}^{8\,\text{km}} C_n^2(h)\,dh = \Lambda(\theta)\,\frac{|F - F_\text{tend}|}{z},
\qquad \Lambda(\theta) = \frac{10^{-12}}{\cos\theta},
\end{equation}
con $\theta$ el ángulo cenital del satélite visto desde el sitio y $z$ el espesor de la capa, 4~km. Seguimos la
convención del artículo: $|F - F_\text{tend}|$ en kelvin y $z$ en kilómetros, con la constante $10^{-12}$ que ellos
ajustaron. Con la integral en la ecuación~\eqref{eq:r0} y el seeing de la ecuación~\eqref{eq:seeing}:
\begin{gather}
r_0 = \left[0{,}423\,\frac{4\pi^2}{\lambda^2}\int C_n^2\,dh\right]^{-3/5},\\
\text{FWHM}_\text{sat} = 0{,}98\,\frac{\lambda}{r_0}
= 0{,}58\,\lambda^{-1/5}\left[4\pi^2\,\Lambda(\theta)\,\frac{|F - F_\text{tend}|}{z}\right]^{3/5}
\quad\text{(radianes, con $\lambda$ en metros)} .
\end{gather}
El 0,58 sale de $0{,}98 \times 0{,}423^{3/5}$. TurbuLENS lo pasa a segundos de arco y calcula también $r_0$. Como el seeing
va como $z^{-3/5}$, cambiar el espesor de la capa cambia todos los valores en la misma proporción: la forma de la serie
no cambia.

\subsection{El ángulo del satélite}

El factor $\Lambda(\theta)$ tiene en cuenta que el satélite ve la columna de aire de lado si el sitio está lejos del punto
bajo él. Para un satélite geoestacionario a una distancia $r = 42\,164$~km del centro de la Tierra, sobre la longitud
$\lambda_s$, visto desde un sitio de latitud $\varphi$ y longitud $\lambda$ (con $R = 6371$~km):
\begin{equation}
\cos\gamma = \cos\varphi\,\cos(\lambda - \lambda_s),\qquad
d = \sqrt{R^2 + r^2 - 2Rr\cos\gamma},\qquad
\cos\theta = \frac{r\cos\gamma - R}{d}.
\end{equation}
Desde el OAN ($4{,}64°$~N, $74{,}08°$~O) y con GOES-Este en $75{,}2°$~O, $\theta = 5{,}62°$ y $\Lambda = 1{,}005\times10^{-12}$:
Bogotá está casi debajo del satélite.

\subsection{El Sol}

De día el Sol calienta el suelo y las capas bajas, y las temperaturas de brillo cambian por razones que no son la
turbulencia nocturna; el método se pensó para la noche. TurbuLENS calcula la altura del Sol en cada instante con las
fórmulas aproximadas del \emph{Almanaque Astronómico} (precisión de medio grado, de sobra para esto) y clasifica cada hora:
\begin{center}\small
\begin{tabular}{@{}l l@{}}
\toprule
Altura del Sol & Estado\\
\midrule
más de 0° & día\\
entre 0° y $-18°$ & crepúsculo\\
menos de $-18°$ & noche (astronómica)\\
\bottomrule
\end{tabular}
\end{center}
Según el criterio elegido, las horas de día (o también las de crepúsculo) quedan fuera de la tendencia y del seeing.

\begin{nota}[Lo que este método puede y no puede decir.]
Mide la capa de 4 a 8~km, no la turbulencia cerca del suelo, que en un sitio urbano o en un techo suele ser la mayor
parte. Las nubes lo invalidan (por eso el canal 13 se muestra al lado). Sus constantes se ajustaron en sitios de
montaña, así que en otro lugar el valor absoluto es incierto. Lo que sí da, con buena confianza, es la \textbf{evolución}
de la noche: qué horas fueron más turbulentas en las capas medias. Comparar esa serie con mediciones del DIMM en las
mismas horas es un buen ejercicio de investigación.
\end{nota}
